\uchapter{Résumé}

Les Systèmes d'Information Décisionnels (SID) classiques reposent sur des architectures
d'entrepôts de données conçues pour des environnements métiers stables et des besoins
analytiques fixes. Face à l'accélération des transformations organisationnelles, à la
prolifération de sources de données hétérogènes et à la maturité analytique croissante
des organisations, ces architectures rigides deviennent un obstacle à l'évolution des
systèmes d'aide à la décision.

Cette recherche aborde le problème de la scalabilité dans les architectures décisionnelles
selon deux axes complémentaires~: la \textit{scalabilité structurelle}, qui concerne la
capacité du système à intégrer de nouvelles sources, à gérer les évolutions de schéma et
à s'adapter à la croissance des volumes sans perturber l'architecture existante~; et la
\textit{scalabilité analytique}, qui reflète la progression d'une BI purement descriptive
vers des capacités prédictives et prescriptives intégrées nativement dans le pipeline
décisionnel.

Pour répondre à ces défis, cette étude s'appuie sur la méthodologie Data Vault~2.0 comme
fondation de la couche de modélisation agile, et explore les patterns d'intégration de
modèles analytiques avancés au sein d'architectures modernes (Lakehouse, Delta
Architecture). La thèse apporte trois contributions conceptuelles~: un cadre formel de
dual-scalabilité (C1), une analyse critique de la composabilité entre Data Vault~2.0 et le
paradigme Lakehouse avec un instrument de critères (C2), et une architecture de référence
en couches avec un contrat d'intégration formel évalué selon six critères (C3).

\bigskip

\noindent\textbf{Mots-clés~:} Architecture scalable, Système d'information décisionnel,
Data Vault~2.0, Entrepôt de données agile, Analytique prédictive, Business Intelligence,
Lakehouse, ETL/ELT, Évolution de schéma, Apprentissage automatique.
